\input{preamble.tex}

\title{Отчет по лабораторной работе № 3. \\ Планирование локальной сети организации}
\author{Сущенко Алина Николаевна \\ НПИбд-01-23}
\date{2026}

\begin{document}

\maketitle
\newpage

\tableofcontents

\newpage
\section{Цель работы}
Познакомиться с принципами планирования локальной сети организации на примере учебного заведения с двумя территориально распределенными площадками.

\newpage
\section{Выполнение работы}

\subsection{Физическая структура сети (L1)}
В логической рабочей области \texttt{Packet Tracer} разместим оборудование в соответствии с иерархической моделью:
\begin{itemize}
    \item \textbf{Уровень ядра (Core):} маршрутизатор Cisco 2811 (\texttt{msk-ansusenko-gw-1}).
    \item \textbf{Уровень распределения (Distribution):} коммутаторы Cisco 2960 (\texttt{msk-ansusenko-sw-1}, \texttt{msk-ansusenko-sw-2}, \texttt{msk-ansusenko-sw-3}, \texttt{msk-ansusenko-sw-4}).
    \item \textbf{Уровень доступа (Access):} коммутатор Cisco 2950 (\texttt{msk-ansusenko-1}) и оконечные устройства (серверы, ПК пользователей).
\end{itemize}
Схема физических подключений с указанием номеров портов представлена на рис. \ref{01}. Номера портов подписаны рядом с каждой линией связи, что является отличительной особенностью схемы физического уровня.

\begin{center}
    \centering
    \includegraphics[width=\textwidth]{../images/01.png}
    \captionof{figure}{Физическая схема сети (L1) с указанием портов подключения.}
    \label{01}
\end{center}

\subsection{Логическая структура сети (L2) и план VLAN}
В соответствии с заданием, сеть разделена на виртуальные локальные сети (VLAN) для изоляции трафика различных групп пользователей и служебных сегментов. План VLAN представлен в табл. \ref{tab:vlan}.

\begin{table}[h]
    \centering
    \caption{План VLAN}
    \label{tab:vlan}
    \begin{tabular}{|c|c|c|}
        \hline
        \textbf{№ VLAN} & \textbf{Имя VLAN} & \textbf{Примечание} \\
        \hline
        1 & default & Не используется \\
        \hline
        2 & management & Для управления устройствами \\
        \hline
        3 & servers & Для серверной фермы \\
        \hline
        4-100 & - & Зарезервировано \\
        \hline
        101 & dk & Дисплейные классы (ДК) \\
        \hline
        102 & departments & Кафедры \\
        \hline
        103 & adm & Администрация \\
        \hline
        104 & other & Для других пользователей \\
        \hline
    \end{tabular}
\end{table}

Схема распределения VLAN по магистральным каналам (trunk) и портам доступа представлена на рис. \ref{02}.

\begin{center}
    \centering
    \includegraphics[width=\textwidth]{../images/02.png}
    \captionof{figure}{Схема VLAN (L2).}
    \label{02}
\end{center}

\subsection{План IP-адресации (L3)}
План IP-адресации разработан для трех частных диапазонов: исходного \texttt{10.128.0.0/16}, а также для \texttt{172.16.0.0/12} и \texttt{192.168.0.0/16}. Структура сетей и VLAN сохранена, изменяются только номера подсетей.

\subsubsection{План для сети 10.128.0.0/16 (исходный)}
Схема маршрутизации для исходного проекта представлена на рис. \ref{03}. Детальный план IP-адресов приведен в табл. \ref{tab:ip1}.

\begin{center}
    \centering
    \includegraphics[width=\textwidth]{../images/03.png}
    \captionof{figure}{Схема маршрутизации (L3) для сети 10.128.0.0/16.}
    \label{03}
\end{center}

\begin{table}[h]
    \centering
    \caption{План IP-адресов для сети 10.128.0.0/16}
    \label{tab:ip1}
    \begin{tabular}{|c|c|c|c|}
        \hline
        \textbf{Название подсети} & \textbf{IP-сеть} & \textbf{VLAN} & \textbf{Адрес шлюза} \\
        \hline
        Серверная ферма & 10.128.0.0/24 & 3 & 10.128.0.1 \\
        Управление & 10.128.1.0/24 & 2 & 10.128.1.1 \\
        Point-to-Point & 10.128.2.0/24 & - & 10.128.2.1 \\
        Дисплейные классы (ДК) & 10.128.3.0/24 & 101 & 10.128.3.1 \\
        Кафедры (К) & 10.128.4.0/24 & 102 & 10.128.4.1 \\
        Администрация (А) & 10.128.5.0/24 & 103 & 10.128.5.1 \\
        Другие пользователи (Д) & 10.128.6.0/24 & 104 & 10.128.6.1 \\
        \hline
    \end{tabular}
\end{table}

\subsubsection{План для сети 172.16.0.0/12}
Для данного диапазона выбрана базовая сеть \texttt{172.16.0.0/16} для сохранения структуры /24 подсетей. План адресации представлен в табл. \ref{tab:ip2}. Физическая схема (L1) и схема VLAN (L2) полностью аналогичны рис. \ref{01} и \ref{02}.

\begin{table}[h]
    \centering
    \caption{План IP-адресов для сети 172.16.0.0/12 (на базе 172.16.0.0/16)}
    \label{tab:ip2}
    \begin{tabular}{|c|c|c|c|}
        \hline
        \textbf{Название подсети} & \textbf{IP-сеть} & \textbf{VLAN} & \textbf{Адрес шлюза} \\
        \hline
        Серверная ферма & 172.16.0.0/24 & 3 & 172.16.0.1 \\
        Управление & 172.16.1.0/24 & 2 & 172.16.1.1 \\
        Point-to-Point & 172.16.2.0/24 & - & 172.16.2.1 \\
        Дисплейные классы (ДК) & 172.16.3.0/24 & 101 & 172.16.3.1 \\
        Кафедры (К) & 172.16.4.0/24 & 102 & 172.16.4.1 \\
        Администрация (А) & 172.16.5.0/24 & 103 & 172.16.5.1 \\
        Другие пользователи (Д) & 172.16.6.0/24 & 104 & 172.16.6.1 \\
        \hline
    \end{tabular}
\end{table}
\begin{center}
    \centering
    \includegraphics[width=\textwidth]{../images/04.png}
    \captionof{figure}{Схема маршрутизации (L4) для сети 172.16.0.0/12.}
    \label{03}
\end{center}
\subsubsection{План для сети 192.168.0.0/16}
Данный диапазон идеально подходит под структуру /16. План адресации представлен в табл. \ref{tab:ip3}.

\begin{table}[h]
    \centering
    \caption{План IP-адресов для сети 192.168.0.0/16}
    \label{tab:ip3}
    \begin{tabular}{|c|c|c|c|}
        \hline
        \textbf{Название подсети} & \textbf{IP-сеть} & \textbf{VLAN} & \textbf{Адрес шлюза} \\
        \hline
        Серверная ферма & 192.168.0.0/24 & 3 & 192.168.0.1 \\
        Управление & 192.168.1.0/24 & 2 & 192.168.1.1 \\
        Point-to-Point & 192.168.2.0/24 & - & 192.168.2.1 \\
        Дисплейные классы (ДК) & 192.168.3.0/24 & 101 & 192.168.3.1 \\
        Кафедры (К) & 192.168.4.0/24 & 102 & 192.168.4.1 \\
        Администрация (А) & 192.168.5.0/24 & 103 & 192.168.5.1 \\
        Другие пользователи (Д) & 192.168.6.0/24 & 104 & 192.168.6.1 \\
        \hline
    \end{tabular}
\end{table}

\subsection{План подключения интерфейсов оборудования}
План распределения портов коммутаторов по VLAN и магистральным каналам является единым для всех трех проектов, так как физическая и логическая топология не меняются. Детальный план представлен в табл. \ref{tab:ports}.

\begin{table}[h]
    \centering
    \caption{План подключения портов оборудования}
    \label{tab:ports}
    \begin{tabular}{|c|c|c|c|c|}
        \hline
        \textbf{Устройство} & \textbf{Порт} & \textbf{Примечание} & \textbf{Access VLAN} & \textbf{Trunk VLAN} \\
        \hline
        \texttt{msk-ansusenko-gw-1} & f0/0 & UpLink & - & 2,3,101,102,103,104 \\
        & f0/1 & к \texttt{msk-ansusenko-sw-1} & - & 2,3,101,102,103,104 \\
        \hline
        \multirow{4}{*}{\texttt{msk-ansusenko-sw-1}} & f0/24 & к \texttt{gw-1} & - & 2,3,101,102,103,104 \\
        & g0/1 & к \texttt{sw-2} & - & 2,3 \\
        & g0/2 & к \texttt{sw-4} & - & 2,101,102,103,104 \\
        & f0/1 & к \texttt{ansusenko-sw-1} & - & 2,101,104 \\
        \hline
        \multirow{2}{*}{\texttt{msk-ansusenko-sw-2}} & g0/1 & к \texttt{sw-1} & - & 2,3 \\
        & g0/2 & к \texttt{sw-3} & - & 2,3 \\
        \hline
        \multirow{5}{*}{\texttt{msk-ansusenko-sw-3}} & g0/1 & к \texttt{sw-2} & - & 2,3 \\
        & f0/1 & Web-сервер & 3 & - \\
        & f0/2 & File-сервер & 3 & - \\
        & f0/3 & Mail-сервер & 3 & - \\
        & f0/4 & DNS-сервер & 3 & - \\
        \hline
        \multirow{5}{*}{\texttt{msk-ansusenko-sw-4}} & g0/1 & к \texttt{sw-1} & - & 2,101,102,103,104 \\
        & f0/1-5 & Пользователи ДК & 101 & - \\
        & f0/6-10 & Пользователи кафедр & 102 & - \\
        & f0/11-15 & Администрация & 103 & - \\
        & f0/16-24 & Другие пользователи & 104 & - \\
        \hline
        \multirow{3}{*}{\texttt{msk-ansusenko-sw-1}} & f0/24 & к \texttt{sw-1} (Донская) & - & 2,101,104 \\
        & f0/1-15 & Пользователи ДК & 101 & - \\
        & f0/16-24 & Другие пользователи & 104 & - \\
        \hline
    \end{tabular}
\end{table}

\newpage
\section{Выводы}
В результате выполнения лабораторной работы были изучены принципы иерархического проектирования локальных сетей.
\begin{enumerate}
    \item Разработана физическая схема сети (L1) для организации с двумя площадками, с четким разделением оборудования по уровням (ядро, распределение, доступ) и указанием номеров портов.
    \item Составлен план VLAN (L2), обеспечивающий логическую сегментацию сети на группы пользователей и служебные сегменты (управление, серверы), что повышает безопасность и управляемость сети.
    \item Разработаны планы IP-адресации (L3) для трех частных диапазонов: \texttt{10.128.0.0/16}, \texttt{172.16.0.0/12} и \texttt{192.168.0.0/16}. Продемонстрирована универсальность разработанной логической структуры: при смене адресного пространства физическая топология и план VLAN остаются неизменными, меняются только IP-адреса подсетей.
    \item Составлен детальный план подключения интерфейсов оборудования, позволяющий однозначно определить конфигурацию каждого порта (access/trunk, принадлежность к VLAN).
\end{enumerate}
Таким образом, поставленные задачи выполнены в полном объеме, полученные навыки могут быть применены при реальном проектировании и развертывании сети предприятия.

\end{document}
